

\section{Variables and Expressions}

\subsection*{Additional: PEMDAS}
\begin{itemize}
    \item Parenthesis, Exponent, Multiplication/Division, Addition/Subtraction
    \item \kr{식의 계산순서를 미국에서 배우는 방법}:
    \begin{itemize}
        \item \kr{괄호 (Parenthesis)}
        \item \kr{제곱 (Exponent)}
        \item \kr{곱셈/나눗셈 (Multiplication/Division)}
        \item \kr{덧셈/뺄셈 (Addition/Subtraction)}
    \end{itemize}
    \item \kr{참고: 미국에서는 소괄호 외의 괄호는 잘 사용하지 않음}
\end{itemize}

\section{Simplifying Expressions}

\subsection*{Additional: Transposition}
\begin{itemize}
    \item \kr{Transposition(이항)은 방정식을 풀 때 중요함}
    \item \kr{이항은 ``=''를 지나갈 때 이항하는 수의 앞에 붙은 부호를 바꾸어줌}
    \item \kr{예시}:
    \begin{align*}
        3x + 7 &= -3x - 5 \\
        7 &= -3x - 3x - 5 \\
        7 &= -6x - 5 \\
        7 + 5 &= -6x \\
        12 &= -6x \\
        x &= -2
    \end{align*}
    \item \kr{곱셈과 나눗셈도 이항 가능}:
    \begin{align*}
        \frac{3x}{7} &= 6 \\
        3x &= 7 \times 6 \\
        3x &= 42 \\
        x &= 14
    \end{align*}
\end{itemize}

\section{Multi-Step Equations}

\subsection*{Additional: Steps Order}
\begin{itemize}
    \item \kr{괄호 풀기}
    \item \kr{동류항 정리}
    \item \kr{이항}
    \item \kr{계수 나누기}
    \item \kr{미국에서는 ``simplify both sides → isolate variable''로 설명}
    \item \kr{한국에서는 ``좌변/우변'' 용어 자주 사용}
\end{itemize}

\section{Inequalities}

\subsection*{Additional: Graphing Inequalities}
\begin{itemize}
    \item \kr{Open circle ($<$, $>$): 포함하지 않음}
    \item \kr{Closed circle ($\leq$, $\geq$): 포함함}
\end{itemize}
